% 强连通分量 = 前向可达集 ∩ 后向可达集。弧：1->2, 2->1, 2->3, 3->4, 4->3, 5->6, 6->4, 6->7, 7->1, 7->5。
% 从 1 前向可达：{1,2,3,4}；能到达 1：{1,2,5,6,7}；交集 {1,2}。
\begin{tikzpicture}[font=\figurefont\fontsize{8}{10}\selectfont,>=Stealth,
  vtx/.style={circle,draw=FigureStroke,line width=.8pt,minimum size=6.2mm,inner sep=0pt,text=FigureInk,fill=white},
  arc/.style={draw=FigureStroke,line width=.8pt,->,shorten >=.5pt},
  root/.style={double,double distance=.9pt},
  note/.style={font=\FigureNoteFont\fontsize{7.5}{9}\selectfont,text=FigureMuted,align=center}]
\def\sccgraph#1#2#3{%
  \node[vtx,fill=#1] (v5) at (0,0) {5};
  \node[vtx,fill=#1] (v7) at (0.95,0.95) {7};
  \node[vtx,fill=#1] (v6) at (0.95,-0.95) {6};
  \node[vtx,root,fill=#2] (v1) at (2.35,0.95) {1};
  \node[vtx,fill=#2] (v2) at (3.55,0.95) {2};
  \node[vtx,fill=#3] (v4) at (2.35,-0.95) {4};
  \node[vtx,fill=#3] (v3) at (3.55,-0.95) {3};
  \draw[arc] (v5)--(v6); \draw[arc] (v6)--(v7); \draw[arc] (v7)--(v5);
  \draw[arc,bend left=18] (v1) to (v2); \draw[arc,bend left=18] (v2) to (v1);
  \draw[arc,bend left=18] (v3) to (v4); \draw[arc,bend left=18] (v4) to (v3);
  \draw[arc] (v7)--(v1); \draw[arc] (v2)--(v3); \draw[arc] (v6)--(v4);}
\begin{scope}
  \sccgraph{white}{FigureGreenFill}{FigureGreenFill}
  \node[note] at (1.8,-1.75) {从$1$出发可达：$\{1,2,3,4\}$};
  \node at (1.8,-2.4) {(a) 前向可达集};
\end{scope}
\begin{scope}[xshift=4.75cm]
  \sccgraph{FigureBlueFill}{FigureBlueFill}{white}
  \node[note] at (1.8,-1.75) {能到达$1$：$\{1,2,5,6,7\}$};
  \node at (1.8,-2.4) {(b) 后向可达集};
\end{scope}
\begin{scope}[xshift=9.5cm]
  \sccgraph{white}{FigureYellowFill}{white}
  \node[note] at (1.8,-1.75) {两者的交：$\{1,2\}$};
  \node at (1.8,-2.4) {(c) $1$所在的强连通分量};
\end{scope}
\end{tikzpicture}
